\documentclass[12pt]{article}
\input{iati_structure.tex}
\renewcommand{\bottomtitlespace}{2cm}

\usepackage[english]{babel}

\begin{document}

\renewcommand{\headerLang}{English}
\renewcommand{\problemName}{AB23. TRIANGLE}

\renewcommand{\headerLeft}{IATI Day 2, Senior group}
\renewcommand{\headerRightFirst}{XVII INTERNATIONAL ADVANCED TOURNAMENT IN INFORMATICS}
\renewcommand{\headerRightSecond}{ONLINE 2026}

\renewcommand{\tl}{$10$ sec.}
\renewcommand{\ml}{$1024$ MB}
\problem{Problem \problemName}
\addauthor{Author: Boris Mihov}{2026}{04}{19}

The new leadership of the national committee, represented by the coordinators of groups A and B, wishes to maintain a strict syllabus of the given problems for selecting the national team. For this purpose, they will need to give a geometry problem. Moreover -- the number of interactive problems given so far are below standard. To fix the crisis in the selection for \texttt{IOI}, Sashka decided to give a problem that is both interactive and geometric. It has the following statement:

The jury has hidden a permutation $P_0,P_1,...,P_{N-1}$ of $\{1,2,3,...,N\}$. You must find the permutation. For this purpose, you may ask the jury the following question: ``Is it possible to form a triangle with positive area with sides $P_A,P_B,P_C$?''

Note that it is known (by the triangle inequality) this is possible if and only if:

\[P_A + P_B > P_C,\ P_A + P_C > P_B \quad \text{and} \quad P_B + P_C > P_A.\]

Write a program \textbf{\texttt{triangle}}, containing a function \texttt{solve}, which will be compiled together with the jury program and will communicate with it by asking questions of the type described above. At the end of its execution, it must determine the permutation.

\subsection{Implementation details}
You should implement the \texttt{solve} function:
\begin{lstlisting}
std::vector<int> solve(int N)
\end{lstlisting}
\begin{itemize}
	\item $N$: length of the permutation.
\end{itemize}

This function will be called $T$ times per test -- once per subtest, each with equal $N$ and it should return the hidden permutation for the subtest. In order to do this, your program can call the jury function \texttt{query}:
\begin{lstlisting}
 bool query(int A, int B, int C)
\end{lstlisting}
\begin{itemize}
	\item $A$, $B$, $C$: indices of the sides $P_A, P_B, P_C$.
\end{itemize}

The function returns \texttt{true}, if a triangle with positive area with sides $P_A,P_B,P_C$ exists and \texttt{false}, otherwise.

\subsection{Constraints}
\vspace{0.1em}
\begin{itemize}
	\item $N = T = 1~000$
\end{itemize}

\subsection{Subtask}
\begin{table}[H]
	\begin{tblr}{width=\linewidth,colspec={|Q[c,m]|Q[c,m]|X[c,m]|}}
		\hline
		\textbf{Subtask} & \textbf{Points} & \textbf{Description}\\
		\hline
		$1$ & $100$ & The subtask consists of $1$ test with $T=1~000$ randomly and uniformly sampled from the set of all possible permutations, each of $N=1~000$ elements.  \\ 
		\hline
	\end{tblr}
	\caption*{The points for a given subtask are obtained only if all the tests for it are successfully passed.}
\end{table}
\FloatBarrier

\newpage
\subsection{Scoring}

Let $Q_{contestant}$ be the average number of queries that your program makes for one call of \texttt{solve} on the single test, and let $Q_{target}=8770$. Then your score for the problem will be:

\[
\begin{cases}
0 & \text{if } Q_{contestant} > 2\times 10^6, \\
100 & \text{if } Q_{contestant} \le Q_{target}, \\
100 \times \max\left(0.15,\; 1-\sqrt{1-\left(\frac{Q_{target}}{Q_{contestant}}\right)^{0.55}}\right) & \text{if } Q_{contestant} > Q_{target}.
\end{cases}
\]

\subsection{Sample interaction}
For this sample interaction $N=4$, $T=2$.
\begin{table}[H]
	\begin{tblr}{|c|c|}
		\hline
		\textbf{Contestant action} & \textbf{Jury action} \\
		\hline
		& \texttt{solve(4)} \\
		\hline
		\texttt{query(0, 0, 0)} & \texttt{true} \\
		\hline
        \texttt{query(0, 1, 2)} & \texttt{false} \\
		\hline
        \texttt{query(0, 1, 3)} & \texttt{false} \\
		\hline
        \texttt{query(0, 2, 3)} & \texttt{true} \\
		\hline
        \texttt{query(1, 2, 3)} & \texttt{false} \\
		\hline
        \texttt{return \{3, 1, 2, 4\};} & \\
        \hline
        & \texttt{solve(4)} \\
		\hline
		\texttt{query(0, 1, 2)} & \texttt{false} \\
		\hline
        \texttt{query(0, 1, 3)} & \texttt{true} \\
		\hline
        \texttt{query(0, 2, 3)} & \texttt{false} \\
		\hline
        \texttt{query(1, 2, 3)} & \texttt{false} \\
		\hline
        \texttt{return \{4, 2, 1, 3\};} & \\
		\hline
	\end{tblr}
\end{table}
\FloatBarrier
	
\subsection{Local testing}
Input format:
\begin{itemize}
	\item line $1$: three integers $T$, $N$, and $R$ -- the number of subtests, the size of the permutations, and the execution mode. If $R = 1$, the local grader will generate uniformly random permutations and will expect a number on the second line -- $S$, which will be the seed for its random generator. If $R = 2$, the input continues as follows:
	\item lines $2$ to $(1+T)$: permutations of $\{1,2,\dots,N\}$.
\end{itemize}

Output format:
\begin{itemize}
	\item line $1$: an error message or the average number of queries for the subtests if all permutations are correctly identified.
\end{itemize}

\end{document}